% Options for packages loaded elsewhere
\PassOptionsToPackage{unicode}{hyperref}
\PassOptionsToPackage{hyphens}{url}
%
\documentclass[
]{article}
\usepackage{amsmath,amssymb}
\usepackage{lmodern}
\usepackage{iftex}
\ifPDFTeX
  \usepackage[T1]{fontenc}
  \usepackage[utf8]{inputenc}
  \usepackage{textcomp} % provide euro and other symbols
\else % if luatex or xetex
  \usepackage{unicode-math}
  \defaultfontfeatures{Scale=MatchLowercase}
  \defaultfontfeatures[\rmfamily]{Ligatures=TeX,Scale=1}
\fi
% Use upquote if available, for straight quotes in verbatim environments
\IfFileExists{upquote.sty}{\usepackage{upquote}}{}
\IfFileExists{microtype.sty}{% use microtype if available
  \usepackage[]{microtype}
  \UseMicrotypeSet[protrusion]{basicmath} % disable protrusion for tt fonts
}{}
\makeatletter
\@ifundefined{KOMAClassName}{% if non-KOMA class
  \IfFileExists{parskip.sty}{%
    \usepackage{parskip}
  }{% else
    \setlength{\parindent}{0pt}
    \setlength{\parskip}{6pt plus 2pt minus 1pt}}
}{% if KOMA class
  \KOMAoptions{parskip=half}}
\makeatother
\usepackage{xcolor}
\setlength{\emergencystretch}{3em} % prevent overfull lines
\providecommand{\tightlist}{%
  \setlength{\itemsep}{0pt}\setlength{\parskip}{0pt}}
\setcounter{secnumdepth}{-\maxdimen} % remove section numbering
\ifLuaTeX
  \usepackage{selnolig}  % disable illegal ligatures
\fi
\IfFileExists{bookmark.sty}{\usepackage{bookmark}}{\usepackage{hyperref}}
\IfFileExists{xurl.sty}{\usepackage{xurl}}{} % add URL line breaks if available
\urlstyle{same} % disable monospaced font for URLs
\hypersetup{
  hidelinks,
  pdfcreator={LaTeX via pandoc}}

\author{}
\date{}

\begin{document}

\documentclass{article}
\usepackage{amsmath}
\usepackage{amsfonts}
\begin{document}

\section*{Lecture Scribe}

\subsection*{CSE400 – Fundamentals of Probability in Computing}

\subsection*{Lecture 21 and 22: Inequalities and Tail Bound}

---

\section*{Outline}

The lecture covers the following topics:

\begin{itemize}
\item Tail and Motivation of Tail Bound
\item Markov’s Inequality
\begin{itemize}
\item Definition and Proof
\item Example
\end{itemize}
\item Chebyshive’s Inequality
\begin{itemize}
\item Definition and Proof
\item Example
\end{itemize}
\item Chernof Bound and Poison Tail
\begin{itemize}
\item Definition and Proof
\item Example
\end{itemize}
\item Jensen’s Inequality
\begin{itemize}
\item Convex Functions
\item Proof and Example
\end{itemize}
\item Case Study and System Level Understanding
\end{itemize}

---

\section*{1. Tails}

\subsection*{Definition}

The tail of a random variable $(X)$ is

[
P{X > x}.
]

\subsection*{Why do we care about tails?}

\begin{itemize}
\item Jobs delayed more than 24 hours
\item Packets dropped due to full buffer
\end{itemize}

\subsection*{Key idea}

\begin{itemize}
\item Mean $(\rightarrow)$ average behavior (easy)
\item Tail $(\rightarrow)$ rare extreme events (hard)
\end{itemize}

\subsection*{Why are tails hard to compute?}

\begin{itemize}
\item Requires summing many small probabilities
\end{itemize}

---

\section*{2. Tails Example-1/2}

\subsection*{Question}

Suppose you’re distributing $(n)$ jobs among $(n)$ servers at random. What’s the probability that a particular server gets $(\geq k)$ jobs?

\subsection*{Statements presented in the lecture}

\begin{itemize}
\item Jobs are assigned randomly
\item Most servers get few jobs
\item Occasionally, one server gets many jobs
\end{itemize}

This is a tail event (rare overload).

\subsection*{Distribution}

[
X \sim \text{Binomial}(n,p)
]

\subsection*{Tail probability}

[
P{X \geq k}
]
[
=

]

[
\sum_{i=k}^{n}
\binom{n}{i}p^i(1-p)^{n-i}
]

The highlighted server has $(\geq k)$ jobs.

---

\section*{3. Tails Example-2/2}

\subsection*{Question}

Jobs arrive to a datacenter according to a Poisson process with rate $(\lambda)$ jobs/hour. What’s the probability that $(\geq k)$ jobs arrive during the first hour?

\subsection*{Statements presented in the lecture}

\begin{itemize}
\item Jobs arrive randomly over time
\item Most hours have moderate arrivals
\item Sometimes, we see a burst of arrivals
\end{itemize}

This is a tail event (sudden spike in load).

\subsection*{Distribution}

[
X \sim \text{Poisson}(\lambda)
]

\subsection*{Tail probability}

[
P{X \geq k}
]
[
=

]

[
\sum_{i=k}^{\infty} e^{-\lambda}\cdot \frac{\lambda^i}{i!}
]

This corresponds to many arrivals within one hour $((\geq k \text{ jobs}))$.

---

\section*{4. Why Tail Bounds?}

Exact tail probabilities are often hard to compute.

\subsection*{Exact computation}

[
P{X \geq k}
]
[
=

]

[
\sum_{i=k}^{\infty} e^{-\lambda}\cdot \frac{\lambda^i}{i!}
]

[
P{X \geq k}
]
[
=

]

[
\sum_{i=k}^{n}
\binom{n}{i}p^i(1-p)^{n-i}
]

Complicated, long sums.

\subsection*{Idea: Instead of exact value}

\begin{itemize}
\item We find an upper bound
\item Easier to compute
\item Still gives a guarantee
\end{itemize}

[
\text{Exact: } P(X \geq k)=?
]

[
\text{Bound: } P(X \geq k)\leq \text{something simple}
]

\subsection*{Key idea}

Approximate safely instead of computing exactly.

---

\section*{5. Running Example (Key Idea)}

The lecture introduces one running example with the following purpose:

\begin{itemize}
\item We will develop progressively better (tighter) tail bounds
\item We test all bounds on the same example
\end{itemize}

\subsection*{Distribution used in the running example}

[
X \sim \text{Binomial}(n,1/2)
]

The mean is indicated as

[
\mu = n/2.
]

The tail region is shown as

[
P(X \geq 3n/4)
]

and is labeled as getting unusually many heads.

---

\section*{6. Running Example: Full Statement}

\subsection*{Experiment}

Flip a fair coin $(n)$ times.

\subsection*{Distribution}

[
X = \text{number of heads}
]

[
X \sim \text{Binomial}\left(n,\tfrac{1}{2}\right)
]

\subsection*{Mean}

[
E[X] = \frac{n}{2}
]

\subsection*{Question}

[
\text{What is } P\left(X \geq \frac{3}{4}n\right)?
]

---

\section*{7. Markov’s Inequality}

\subsection*{Key idea}

\begin{itemize}
\item Uses only the mean
\item Large values $(\Rightarrow)$ rare
\end{itemize}

\subsection*{Example}

[
\text{Avg} = 10 \Rightarrow \text{few values can be } 100+
]

\subsection*{Statement}

[
P(X \geq a)\leq \frac{E[X]}{a}
]

---

\section*{8. Proof (Idea) of Markov’s Inequality}

The lecture presents the proof idea in the following steps:

\begin{itemize}
\item Consider only $(X \geq a)$
\item Replace $(X)$ by $(a)$
\end{itemize}

Then the derivation shown is

[
E[X]=\sum_{x=0}^{\infty} x\cdot p_X(x)
]

[
\geq \sum_{x=a}^{\infty} x\cdot p_X(x)
]

[
\geq \sum_{x=a}^{\infty} a\cdot p_X(x)
]

[
= a\cdot P(X\geq a)
]

Hence,

[
P(X\geq a)\leq \frac{E[X]}{a}.
]

---

\section*{9. Markov’s Inequality on Running Example}

\subsection*{Question}

Flip a fair coin $(n)$ times:

What’s a tail bound on the probability of getting at least $(\frac{3}{4}n)$ heads?

\subsection*{Given}

[
X = \text{Number Heads} \sim \text{Binomial}\left(n,\frac{1}{2}\right)
]

[
\mu = E[X] = \frac{n}{2}
]

\subsection*{Applying Markov’s inequality}

[
P\left{X \geq \frac{3}{4}n\right}
\leq
\frac{\mu}{\frac{3}{4}n}
]
[
=

]

[
# \frac{\frac{n}{2}}{\frac{3}{4}n}
]

[
\frac{2}{3}
]

\subsection*{Conclusion stated in the lecture}

Clearly a terrible bound because doesn’t involve $(n)$.

---

\section*{10. Markov Inequality: Real-Life Example}

\subsection*{Scenario}

Income Distribution

\begin{itemize}
\item Average income $(= 10{,}000)$
\end{itemize}

---

\section*{11. Case Study}

The lecture states:

How to predict the worst case demands on server?

---

\section*{12. Closing}

Thank You :)

\end{document}

\end{document}
